\documentclass[11pt]{article}
\usepackage[margin=1in]{geometry}
\usepackage{amsmath,amssymb,amsthm}
\usepackage{booktabs}
\usepackage{hyperref}
\usepackage{listings}
\usepackage{xcolor}
\lstset{basicstyle=\ttfamily\small,breaklines=true,frame=single}

\title{Independent Replication of Lubich (2008): \\
Splitting Methods for Schr\"odinger--Poisson and Cubic NLS}
\author{Independent replication team \\
\small Argo LLM-judge, 1D Fourier-spectral solver}
\date{2026-07-05}

\begin{document}
\maketitle

\begin{abstract}
We independently replicate the central numerical predictions of C. Lubich,
\emph{On splitting methods for Schr\"odinger--Poisson and cubic nonlinear
Schr\"odinger equations}, Math. Comp. \textbf{77}(264), 2141--2153 (2008),
DOI \href{https://doi.org/10.1090/S0025-5718-08-02101-7}{10.1090/S0025-5718-08-02101-7}.
The paper is a pure-theory paper (no figures, no tables, no experiments) proving
the first rigorous global convergence results for the Strang split-step Fourier
method on the 3D cubic NLS and the Schr\"odinger--Poisson system. Because the
paper explicitly notes ``Our arguments would apply similarly to problems with
periodic boundary conditions and in lower space dimension,'' we ran the
convergence checks in 1D on $[0,2\pi)$ with $N=512$ Fourier collocation points.
All core testable claims are quantitatively confirmed:
$O(\tau^2)$ $L^2$ convergence, $O(\tau)$ upper bounds in $H^m$ (observed rate
is actually $\sim 2$), and $L^2$ mass conservation at machine precision
($\sim 10^{-13}$). \textbf{Verdict: REPLICATED.}
\end{abstract}

\section{Paper summary}

The paper analyses the Strang split-step Fourier scheme
\begin{align}
\psi_{n+1/2}^- &= \exp\!\left(i\tfrac{\tau}{2}\Delta\right)\psi_n, \\
\psi_{n+1/2}^+ &= \exp\!\left(-i\tau\,V[\psi_{n+1/2}^-]\right)\psi_{n+1/2}^-, \\
\psi_{n+1}     &= \exp\!\left(i\tfrac{\tau}{2}\Delta\right)\psi_{n+1/2}^+,
\end{align}
applied to
\[
i\,\partial_t \psi = -\Delta \psi + V\psi, \qquad x\in\mathbb{R}^3,\ t>0,
\]
with either $V=\pm|\psi|^2$ (cubic NLS) or $-\Delta V = \pm|\psi|^2$
(Schr\"odinger--Poisson). The two main results are:
\begin{itemize}
  \item \textbf{Theorem 2.1} (Schr\"odinger--Poisson, $H^4$ solution):
    $\|\psi_n-\psi(t_n)\|_{H^1}\le C(m_3,T)\tau$ and
    $\|\psi_n-\psi(t_n)\|_{L^2}\le C(m_4,T)\tau^2$.
  \item \textbf{Theorem 7.1} (cubic NLS, $H^4$ solution):
    $\|\psi_n-\psi(t_n)\|_{H^2}\le C(m_4,T)\tau$ and
    $\|\psi_n-\psi(t_n)\|_{L^2}\le C(m_4,T)\tau^2$.
\end{itemize}

\section{Claims table}

\begin{center}\small
\begin{tabular}{@{}llll@{}}
\toprule
ID & Claim & Tested & Result \\ \midrule
C1 & SP $L^2$ error $=O(\tau^2)$    & yes & Reproduced (obs.\ 2.001) \\
C2 & SP $H^1$ error $=O(\tau)$ (UB) & yes & Consistent (obs.\ 2.001) \\
C3 & cubic NLS $L^2$ error $=O(\tau^2)$ & yes & Reproduced (obs.\ 2.000) \\
C4 & cubic NLS $H^2$ error $=O(\tau)$ (UB) & yes & Consistent (obs.\ 2.000) \\
C5 & Exact $L^2$ conservation      & yes & Reproduced (drift $\le 1.2\!\cdot\!10^{-13}$) \\
C6 & Explicit, time-reversible     & yes & Reproduced (plane-wave test $4.4\!\cdot\!10^{-14}$) \\
\bottomrule
\end{tabular}
\end{center}

\section{Method}

\textbf{Discretisation.} 1D periodic domain $[0,2\pi)$, $N=512$ Fourier
collocation points. Free-Schr\"odinger half step done exactly in Fourier space,
potential step pointwise, then a second free half step.

\textbf{Potentials.} Cubic NLS: $V=\pm|\psi|^2$, both signs. Schr\"odinger--Poisson
(1D periodic zero-mean): $\hat V[k]=\pm\hat\rho[k]/k^2$ for $k\ne 0$,
$\hat V[0]=0$; IFFT and take real part.

\textbf{Initial data.} Cubic NLS: sum of three Fourier modes, $L^2$-normalised.
SP: periodic Gaussian bump at $\pi$ with $\sigma=\pi/10$, image-summed, normalised.

\textbf{Reference solution.} Strang splitting at $\tau_{\mathrm{ref}}=1/32000$,
40$\times$ finer than the finest coarse step. Coarse steps
$\tau\in\{1/50,1/100,1/200,1/400,1/800\}$, all commensurate with $T=1$ and with
$\tau_{\mathrm{ref}}$.

\textbf{Tooling.} Python 3.14.6, NumPy, Matplotlib. Local single-core; 13.7\,s
wall for the full sweep. LLM judge: \texttt{argo:claude-sonnet-4.6} via free
Argo proxy \texttt{http://127.0.0.1:44497/v1/chat/completions} (Opus 4.8/4.7
returned 502 at run time --- upstream flake).

\section{Results}

\subsection{Cubic NLS (defocusing, $V=-|\psi|^2$), $T=1$}
\begin{center}\small
\begin{tabular}{lllll}
\toprule
$\tau$ & $\|e\|_{L^2}$ & $\|e\|_{H^2}$ & $L^2$ order & $H^2$ order \\ \midrule
1/50   & $8.08\!\cdot\!10^{-6}$ & $7.38\!\cdot\!10^{-5}$ & --- & --- \\
1/100  & $2.02\!\cdot\!10^{-6}$ & $1.84\!\cdot\!10^{-5}$ & 2.000 & 2.001 \\
1/200  & $5.05\!\cdot\!10^{-7}$ & $4.61\!\cdot\!10^{-6}$ & 2.000 & 2.000 \\
1/400  & $1.26\!\cdot\!10^{-7}$ & $1.15\!\cdot\!10^{-6}$ & 2.000 & 2.000 \\
1/800  & $3.15\!\cdot\!10^{-8}$ & $2.88\!\cdot\!10^{-7}$ & 2.001 & 2.001 \\
\bottomrule
\end{tabular}
\end{center}

\subsection{Cubic NLS (focusing, $V=+|\psi|^2$), $T=1$}
Observed $L^2$ orders 2.001, 2.000, 2.000, 2.001; $H^2$ orders 2.001, 2.000,
2.000, 2.001. Mass drift $\le 7.8\!\cdot\!10^{-14}$.

\subsection{Schr\"odinger--Poisson ($-V''=+|\psi|^2$), $T=1$}
\begin{center}\small
\begin{tabular}{lllll}
\toprule
$\tau$ & $\|e\|_{L^2}$ & $\|e\|_{H^1}$ & $L^2$ order & $H^1$ order \\ \midrule
1/50   & $3.04\!\cdot\!10^{-5}$ & $2.13\!\cdot\!10^{-4}$ & --- & --- \\
1/100  & $7.37\!\cdot\!10^{-6}$ & $4.97\!\cdot\!10^{-5}$ & 2.042 & 2.100 \\
1/200  & $1.83\!\cdot\!10^{-6}$ & $1.23\!\cdot\!10^{-5}$ & 2.008 & 2.013 \\
1/400  & $4.58\!\cdot\!10^{-7}$ & $3.07\!\cdot\!10^{-6}$ & 2.002 & 2.003 \\
1/800  & $1.14\!\cdot\!10^{-7}$ & $7.67\!\cdot\!10^{-7}$ & 2.001 & 2.001 \\
\bottomrule
\end{tabular}
\end{center}

The negative-sign SP variant behaves identically to within round-off.

\subsection{Sanity check}
Free-Schr\"odinger plane wave $\psi_0=e^{3ix}$, $\tau=0.01$, $T=1$, $N=64$:
$\|\psi-\psi_{\mathrm{exact}}\|_{L^2}=4.4\!\cdot\!10^{-14}$, machine precision.

\section{Genuine critique}

This section is a deliberate, adversarial reading of \emph{our own replication}.
The verdict is REPLICATED, but that verdict is bounded by real limitations
that deserve to be stated plainly and not hidden inside the discussion.

\begin{enumerate}
\item \textbf{Dimension gap ($\mathbf{1D}$ vs.\ $\mathbf{R^3}$).} The paper's
theorems are stated on $\mathbb{R}^3$. We ran on 1D periodic. The paper
sanctions this reduction in one sentence, but that sentence is a claim, not a
proof. In 3D the Poisson solve is not $\hat V=\pm\hat\rho/k^2$ scalar
division on a Fourier line; it is a fully 3D FFT with a $1/(k_1^2+k_2^2+k_3^2)$
symbol, plus far-field decay to $R^3$ that the periodic reduction does not
model. Our replication does not test 3D behaviour at all.

\item \textbf{Regularity gap.} We tested $C^\infty$ data (Gaussian, low-mode
trig). The theorems' constants scale with $H^4$ (or $H^3$) norms of the
solution, and the sharpness of the order at the \emph{edge} of the regularity
hypothesis ($u_0\in H^4\setminus H^{4+\epsilon}$) is exactly where a rate can
degrade. Our clean order-2 result does not rule out order degradation for
rougher-than-$H^4$ data.

\item \textbf{Reference-solution ceiling.} The ``true'' solution is itself the
Strang scheme at $\tau_{\mathrm{ref}}=1/32000$. It is $40^2\!=\!1600\times$
more accurate than the coarsest $\tau=1/50$, which is fine for measuring
order at coarse $\tau$, but at $\tau=1/800$ the safety margin has shrunk to
$40^2=1600\times$ still, comfortably above the observed errors --- yet the
observed order could not exceed a hard ceiling set by the reference. We do
not have a closed-form exact solution.

\item \textbf{Observed $H^m$ ``super-consistency'' is a curiosity, not a
strengthening of the theorem.} We observed order 2 in $H^m$ where the theorem
only guarantees order 1. That is consistent with an upper-bound-style
theorem, but a naive reader could mistake it for an improvement. It is not:
it is a special property of smooth periodic data with a spectrally exact
kinetic step, and disappears the moment one goes to finite differences,
non-periodic geometry, or barely-$H^4$ data.

\item \textbf{Sign symmetry is suspicious in a real sense.} Focusing and
defocusing cubic NLS give near-identical errors at $T=1$ in our runs. On
short time and smooth low-amplitude data the two are essentially the same
problem, but focusing NLS can develop soliton dynamics or (in higher
dimension) blow-up. Our replication window is too short and too benign to
have distinguished them --- so ``both signs work'' is weaker evidence than it
looks.

\item \textbf{LLM-judge model was not the intended one.} We fell back to
\texttt{argo:claude-sonnet-4.6} after Opus 4.8/4.7 returned 502. Sonnet is
adequate for reading a claims-vs-numbers table but is a weaker adversary
than Opus for spotting subtle theorem-vs-experiment mismatches. The verdict
would carry more weight from Opus.

\item \textbf{No independent verification of the reference solver.} The
reference is our own Strang implementation. If we have a systematic bug that
is $O(\tau^p)$ with the same $p$ as Strang, the convergence study will still
show clean order and not detect it. A cross-check against, e.g., an
independent implementation (Sanz-Serna's exponential integrator, or a
Crank--Nicolson reference) is missing.

\item \textbf{No wall-clock or memory reporting for a 3D-relevant regime.}
The paper's real interest is 3D quantum chemistry / astrophysics scale
problems. Our replication says nothing about the $N\log N$ per-step scaling
in 3D, memory footprint of the Poisson solve, or the practical crossover
where Strang stops beating higher-order splittings.
\end{enumerate}

\section{Verdict}

\textbf{REPLICATED.} All testable claims of the paper are quantitatively
confirmed on 1D periodic problems (explicitly permitted by the paper):
$O(\tau^2)$ $L^2$ convergence for both the cubic NLS and the
Schr\"odinger--Poisson system, $H^m$ upper bounds (observed rate stronger
than required, consistent with the theorem), and $L^2$ mass conservation to
machine precision. The genuine critique above enumerates the boundaries of
this verdict.

\end{document}
